\documentclass[palette=basic,mode=light]{impaTeXPoster}

\usepackage[utf8]{inputenc}
\usepackage[T1]{fontenc}
\usepackage{lipsum}

\title{Poster Teste da Silva}
\author{Bruno Pereira, Igor Zwirtes, Marcos Melo}
\authoremail{al.bruno.paula@impatech.edu.br, al.igor.zwirtes@impatech.edu.br, al.marcos.melo@impatech.edu.br}
\posterlogo{images/impatech(big).png}
\leftposterlogo{images/tex.png}

\begin{document}
\maketitle

\begin{columns}

% ==========================================
% COLUNA 1
% ==========================================
\column{0.4}

\begin{postercard}[beige,white,black,black]{Introdução}
Este é um exemplo de seção usando a classe ImpaTeXPoster.

\lipsum[1]
\lipsum[1]
\end{postercard}

\plainblock{
    \posterfigure{0.9}{images/lora.png}%
        {Arquitetura utilizada para o treinamento com LoRA.}{fig:lora}
}

\begin{postercard}[beige,white,black,black]{Objetivo}
Mostrar como usar diferentes cores no ambiente.

\lipsum[2]
\end{postercard}

\begin{postercard}[beige,white,black,black]{Metodologia}
Mostrar como usar diferentes cores no ambiente.

\lipsum[2]
\end{postercard}

% ==========================================
% COLUNA 2
% ==========================================
\column{0.4}

\begin{postercard}[beige,white,black,black]{Experimentos}
Este é um exemplo de seção usando a classe ImpaTeXPoster.

\lipsum[1]
\end{postercard}

\plainblock{
    \posterfigure{0.9}{images/llm.jpg}%
        {Pipeline geral de \textit{text-to-SQL}.}{fig:llm}
}

\begin{postercard}[beige,white,black,black]{Resultados}
Mostrar como usar diferentes cores no ambiente.

\lipsum[2]
\lipsum[1]
\lipsum[3]
\end{postercard}

% ==========================================
% COLUNA 3
% ==========================================
\column{0.2}

\begin{postercard}[beige,white,black,black]{Conclusão}
\lipsum[4]
\end{postercard}

\begin{postercard}[beige,white,black,black]{Trabalhos futuros}
\lipsum[3]
\end{postercard}

% Referências no corpo: ordem alfabética (ABNT NBR 6023:2018)
\begin{postercard}[beige,white,black,black]{Referências}
\begin{thebibliography}{99}

\bibitem{chalmers}
CHALMERS, A. F\@.
\textit{O que é ciência, afinal?}
São Paulo: Brasiliense, 1999.

\bibitem{popper}
POPPER, K\@.
\textit{A lógica da pesquisa científica}.
São Paulo: Cultrix, 2002.

\end{thebibliography}
\end{postercard}

\end{columns}

% Rodapé: QR code + referências completas (ABNT NBR 6023:2018)
\posterfooter{https://github.com/impaTeX/modelo-cartaz/}{%
    \begin{thebibliography}{99}

    \bibitem{lora}
    HU, E.~J\@. et al\@.
    \textit{LoRA: Low-rank adaptation of large language models}.
    arXiv preprint arXiv:2106.09685, 2021.

    \bibitem{mimiciii}
    JOHNSON, A.~E.~W\@. et al\@.
    \textit{MIMIC-III, a freely accessible critical care database}.
    v.~3, n.~160035, 2016.

    \bibitem{mimicsql}
    WANG, P\@.; SHI, T\@.; REDDY, C.~K\@.
    \textit{Text-to-SQL generation for question answering on electronic medical records}.
    New York: ACM, 2020. p.~350--361.

    \bibitem{qwen25}
    YANG, A\@. et al\@.
    \textit{Qwen2.5 technical report}.
    arXiv preprint arXiv:2412.15115, 2024.

    \bibitem{spider}
    YU, T\@. et al\@.
    \textit{Spider: A large-scale human-labeled dataset for complex and cross-domain semantic parsing and text-to-SQL task}.
    Brussels: Association for Computational Linguistics, 2018.

    \end{thebibliography}
}

\end{document}